\section*{第18章：案例研究——命令式对象}
\addcontentsline{toc}{section}{第18章：案例研究——命令式对象}

\begin{taplsolutionentrybox}
\phantomsection
\label{sol:translator:18.6.2}
\textbf{练习18.6.2}

在项的语法中增加记录扩展构造：
\[
 t ::= \cdots \mid
 t\ \mathsf{with}\ \{l_i=t_i\}^{i\in1..n}
\]
右侧列表中的标签互不重复；若某个标签已经出现在基记录中，新字段会替换旧
字段，否则追加为新字段。

求值采用从左到右的按值次序。先求值基记录；基记录成为值后，再依次求值扩展
列表中的各字段。所有部分都是值时执行合并：
\[
\inferrule*[right=\textsc{E-WithBase}]
 {t\trans t'}
 {t\ \mathsf{with}\ F\trans t'\ \mathsf{with}\ F}
\]
\[
\scalebox{0.82}{$
\inferrule*[right=\textsc{E-WithField}]
 {t_i\trans t_i'}
 {v\ \mathsf{with}\
  \{l_1=v_1,\ldots,l_i=t_i,\ldots,l_n=t_n\}
  \trans
  v\ \mathsf{with}\
  \{l_1=v_1,\ldots,l_i=t_i',\ldots,l_n=t_n\}}
$}
\]
其中最后一条规则只在 \(v,v_1,\ldots,v_{i-1}\) 已经是值时使用。若
\(v=\{k_j=w_j\}^{j\in1..m}\)，合并规则为
\[
\scalebox{0.88}{$
\inferrule*[right=\textsc{E-With}]{ }
 {\{k_j=w_j\}^{j\in1..m}\ \mathsf{with}\
  \{l_i=v_i\}^{i\in1..n}
  \trans
  \{k_j=w_j\mid k_j\notin\{l_1,\ldots,l_n\}\},
  \{l_i=v_i\}^{i\in1..n}}
$}
\]
结论右侧表示先保留未被覆盖的旧字段，再放入全部新字段。

类型规则先推导基记录与每个新字段的类型，再以同样方式合并字段类型：
\[
\inferrule*[right=\textsc{T-With}]
 {\Gamma\vdash t:\{k_j:T_j\}^{j\in1..m}\\
  \Gamma\vdash t_i:S_i\quad(i\in1..n)}
 {\Gamma\vdash
  t\ \mathsf{with}\ \{l_i=t_i\}^{i\in1..n}:
  \{k_j:T_j\mid k_j\notin\{l_1,\ldots,l_n\}\},
    \{l_i:S_i\}^{i\in1..n}}
\]
它允许覆盖字段时改变字段类型，因为扩展结果是一条新记录，不会修改原记录。
若希望构造只表达“增加字段而不覆盖”，再加旁条件
\(\{l_i\}\cap\{k_j\}=\varnothing\) 即可。正文中的例子由此写成
\(\mathit{super}\ \mathsf{with}\ \{\mathit{reset}=\cdots\}\)，无需重复
列出 \texttt{get} 与 \texttt{inc}。

\taplsolutionbacklink{ex:18.6.2}{返回练习18.6.2}
\end{taplsolutionentrybox}
